\beginsong{Regarde l'étoile}[by={C. Blanchard / d'après Saint Bernard de Clairvaux}]
\beginchorus
Re\[Am]garde l'étoile, in\[F]voque Marie,
\[Dm]Si tu la \[Am]suis, tu ne crains \[E]rien !
Re\[Am]garde l'étoile, inv\[F]oque Marie,
\[Dm]Elle te con\[Am]duit sur le che\[E]min !
\endchorus
\beginverse
Si le \[Am]vent des tentations s'é\[F]lève,
Si tu \[Dm]heurtes le rocher des é\[E]preuves,
Si les \[Am]flots de l'ambition t'en\[F]traînent,
Si l'o\[Dm]rage des passions se dé\[E]chaîne :
\endverse
\beginverse
Dans l'an\[Am]goisse et les périls, le \[F]doute,
Quand la \[Dm]nuit du désespoir te re\[E]couvre,
Si de\[Am]vant la gravité de tes \[F]fautes,
La pen\[Dm]sée du jugement te tour\[E]mente :
\endverse
\beginverse
Si ton \[Am]âme est envahie de co\[F]lère,
Jalou\[Dm]sie et trahison te sub\[E]mergent,
Si ton \[Am]c{\oe}ur est englouti dans le \[F]gouffre,
Empor\[Dm]té par les courants de tris\[E]tesse :
\endverse
\beginverse
Elle se \[Am]lève sur la mer, elle é\[F]claire,
Son é\[Dm]clat et ses rayons illu\[E]minent,
Sa lu\[Am]mière resplendit sur la \[F]terre,
Dans les \[Dm]cieux et jusqu'au fond des a\[E]bîmes.
\endverse
\beginverse*
\[Am]Si tu la suis, tu ne dé\[F]vies pas,
\[Dm]Si tu la pries, tu ne fai\[E]blis pas,
\[Am]Tu ne crains rien, elle est a\[F]vec toi,
\[Dm]Et jusqu'au port, elle te guide\[E]ra.
\endverse
\endsong
